
\begin{abstract}
长江十年禁渔实施后，四大家鱼资源量回升，长江江豚和长江鲟种群亦呈恢复态势；资源枯竭、江湖连通性受阻、水体污染与外来物种入侵却仍同时存在。为回答"资源恢复能否稳定转化为生态系统健康改善"，本文把流域平均食物网、珍稀物种响应、Hastings-Powell三级食物链和综合情景评价置于同一分层动力学模型中，考察禁渔收益从基础资源层向经济鱼类、珍稀物种及长期动力学行为传递时发生的变化。

问题一按食性保留四类基础资源与四大家鱼功能群的配对关系，用2025年四大家鱼资源量恢复至禁渔前1.8倍标定唯一增益参数，湖北江段单网次渔获量提升120\%作为局地对照。结果显示：2025年四大家鱼总生物量为禁渔前的1.805倍，单位捕捞努力量渔获量($\mathrm{CPUE}^*$)为1.838倍；摄食压力指数升至4.213时，基础资源量末值降至0.204。鱼群恢复与基础资源承压由此同时出现。

问题二接收问题一的资源--鱼群轨迹，以综合生态状态指数调节中层猎物承载力，建立珍稀物种种群响应模型，并设置高阻隔、无人工放流和污染三种对照情景。基线情景下，长江江豚与长江鲟种群末值为禁渔前的1.133倍和1.140倍；高阻隔情景降至0.901倍和0.856倍，无人工放流时长江鲟降至0.519倍，污染情景下二者为0.610倍和0.704倍。结果显示，通道阻隔同时压低两类物种，放流变化则主要作用于长江鲟。

问题三、四采用统一的Hastings-Powell三级食物链模型。承载力参数$K$增大时，轨道由规则有界振荡转为复杂非线性波动；代表点$K=1.795$、$2.665$、$4.800$的峰谷差和峰间隔离散程度均增大。短时间窗分析显示，首个连续正李雅普诺夫指数区间位于$K=3.2186$附近，但75组扩展测试中仅5组在$K=3.20\sim3.22$区间保留由负转正。故本文仅保留"系统动力学复杂性随承载力提高而增强"的方向性结论。

问题五在同一积分后期时间窗内比较饵料资源、长江江豚、长江鲟和入侵压力，采用熵权-CRITIC组合赋权模型，取二者平均后得到三种情景得分0.662、0.423和0.155。组合权重经2000次局部扰动，情景排序稳定率仍为100\%。这一递减次序对应"水体污染首先削弱禁渔收益，外来物种入侵进一步导致系统状态恶化"。

十年禁渔已经带来资源恢复，但资源量回升不会自然转化为系统整体健康改善。若不同步推进水质治理、江湖连通性修复、珍稀物种定向放流与入侵防控，禁渔收益在向高营养级传递及多重压力叠加过程中仍可能逐步减弱。
\end{abstract}

